\chapter{Evaluation}
\label{ch:evaluation}

% =====================================================================
% Reusable coverage-over-time plot.
%   \covtimeseries{<sut>}{<suffix>}{<xcol>}{<ycol>}{<xlabel>}{<ylabel>}
% Reads data/<sut>/synth_<suffix>.csv and data/<sut>/direct_<suffix>.csv,
% plotting whichever exist. The chapter therefore compiles before all data
% is collected, and new SUT/tool runs appear automatically once their CSVs
% are produced by data/extract_results.py - no edits to this file needed.
% =====================================================================
% Optional first arg (#1) injects extra axis options, e.g. [xmax=0.4] to
% truncate a panel; the six positional args are <sut> <suffix> <xcol> <ycol>
% <xlabel> <ylabel> as before.
\newif\ifhascsv
\newcommand{\covtimeseries}[7][]{%
  \hascsvfalse
  \IfFileExists{data/#2/synth_#3.csv}{\hascsvtrue}{}%
  \IfFileExists{data/#2/direct_#3.csv}{\hascsvtrue}{}%
  \ifhascsv
    \begin{tikzpicture}
      \begin{axis}[evalplot, xlabel={#6}, ylabel={#7}, #1]
        \IfFileExists{data/#2/synth_#3.csv}{%
          \addplot table [col sep=comma, x=#4, y=#5] {data/#2/synth_#3.csv};
          \addlegendentry{ELFuzz}}{}%
        \IfFileExists{data/#2/direct_#3.csv}{%
          \addplot table [col sep=comma, x=#4, y=#5] {data/#2/direct_#3.csv};
          \addlegendentry{\toolname{}}}{}%
      \end{axis}
    \end{tikzpicture}%
  \else
    \fbox{\parbox{0.7\linewidth}{\centering\vspace{2cm}\textit{Data pending for #2.}\vspace{2cm}}}%
  \fi
}

% Direct-only coverage-over-time plot (single series). Used for the pool
% evolution figure, where an ELFuzz/\toolname{} comparison is not meaningful
% (ELFuzz's per-generation coverage excludes its large post-evolution
% `produce' step).
\newcommand{\directcovtimeseries}[7][]{%
  \IfFileExists{data/#2/direct_#3.csv}{%
    \begin{tikzpicture}
      \begin{axis}[evalplot, xlabel={#6}, ylabel={#7}, #1]
        \addplot table [col sep=comma, x=#4, y=#5] {data/#2/direct_#3.csv};
      \end{axis}
    \end{tikzpicture}%
  }{%
    \fbox{\parbox{0.7\linewidth}{\centering\vspace{2cm}\textit{Data pending for #2.}\vspace{2cm}}}%
  }%
}

The aim of this project is to test the hypothesis laid out in the introduction (Chapter~\ref{sec:intro}). ELFuzz produces seed inputs \textit{indirectly}: it uses an LLM to synthesise generation-based fuzzers, and those fuzzers are then run to emit the inputs \cite{chen_2025_elfuzz}. \toolname{} removes the intermediate program and asks the LLM for the inputs directly. Our hypothesis is that this is a more direct route to diverse, semantically interesting inputs, and therefore to more varied coverage of the SUT - which is ultimately what a fuzzing campaign is trying to achieve. By eliminating the generator, the model's knowledge of the input format is spent on the inputs themselves rather than on writing programs able to produce them, and coverage feedback selects directly over the objects we care about.

This chapter sets out how we tested that hypothesis. The whole evaluation is a direct comparison between original ELFuzz and \toolname{}, run under identical models, hardware, and time budgets.

\section{Aims and Research Questions}
\label{sec:eval-rqs}

We split evaluation into two research questions.

\begin{itemize}
    \item \textbf{RQ1 (seed quality).} How much SUT coverage does each tool's seed corpus achieve \textit{on its own}, before any mutation fuzzing? This measures the raw quality of what each approach produces, and is the most direct test of the hypothesis: if directly generated inputs are more varied, they should cover more of the SUT.
    \item \textbf{RQ2 (downstream fuzzing).} When the seed corpus is used to bootstrap a grey-box mutation fuzzer (AFL++ \cite{fioraldi_2020_aflplusplus}), how does coverage grow over the campaign, and where does it end? A more diverse corpus should let the fuzzer past shallow checks sooner and reach a higher plateau.
\end{itemize}

\section{Experimental Setup}
\label{sec:eval-setup}

\subsection{Systems under test}

We evaluate on a subset of the SUTs used by ELFuzz, spanning structured data formats and language interpreters \cite{chen_2025_elfuzz}. Table~\ref{tab:suts} lists the two SUTs evaluated so far with their input format and the size of their AFL++ coverage map. We deliberately chose targets at opposite ends of the complexity range: jsoncpp, a small JSON parser, and cpython3, the full CPython interpreter, which is roughly $25\times$ larger by instrumented edges.

\begin{table}[h]
    \centering
    \begin{tabular}{llr}
        \toprule
        \textbf{SUT} & \textbf{Input format} & \textbf{Instrumented edges} \\
        \midrule
        jsoncpp   & JSON   & 5{,}236 \\
        cpython3  & Python & 133{,}395 \\
        \bottomrule
    \end{tabular}
    \caption{The systems under test evaluated so far (a subset of the seven ELFuzz SUTs). ``Instrumented edges'' is the size of the AFL++ coverage map (the total reachable edges of the target), used here as a direct measure of complexity.}
    \label{tab:suts}
\end{table}

\subsection{Configurations compared}

The two configurations are ELFuzz and \toolname{}. To isolate the effect of the design difference, every other factor is held fixed:

\begin{itemize}
    \item \textbf{LLM:} Qwen2.5-Coder-1.5B \cite{hui_2024_qwen25coder}, a small local model chosen to fit the available GPU. Both tools use the same model for their evolution loop.
    \item \textbf{Hardware:} Intel i9-13900K CPU, 32GB RAM, NVIDIA RTX 4070 Ti GPU.
    \item \textbf{Fuzzer:} Both pipelines end by feeding their seed corpus to the same AFL++ configuration.
\end{itemize}

\subsection{Testing Procedure}
\label{sec:eval-protocol}

Each SUT is run end-to-end through the same pipeline for both tools:

\begin{enumerate}
    \item \textbf{Synthesis / evolution (6\,h).} ELFuzz evolves a population of generators; \toolname{} evolves the input pool of Chapter~\ref{ch:project}. Both are given a 6-hour wall-clock budget.
    \item \textbf{Corpus production.} For ELFuzz, a separate \texttt{produce} step runs the evolved elite fuzzers for \textbf{200\,seconds in total}, which is already sufficient to emit between $\sim$0.9 and $\sim$2.8 million candidate inputs. \toolname{} has no \texttt{produce} step, as its evolved pool \textit{is} the corpus. Both corpora are then minimised to remove inputs that are redundant for coverage.
    \item \textbf{Minimisation.} The raw corpus is minimised with \texttt{afl-cmin} to a coverage-preserving subset, a necessary step before seeding the AFL++ campaign.
    \item \textbf{Seed-coverage benchmark (RQ1).} The minimised corpus is executed against the SUT and the total coverage it reaches is recorded.
    \item \textbf{AFL++ campaign (6\,h, RQ2).} The minimised corpus seeds a 6-hour AFL++ campaign, and coverage is logged over time.
\end{enumerate}

\section{Results}
\label{sec:eval-results}

\subsection{Corpus size}
\label{sec:eval-seedcount}

First we observe how large each tool's seed corpus is, before and after minimisation. This is a basic sanity check and also gives context for the seed-coverage results of RQ1, since a larger corpus has more chances to cover more edges.

ELFuzz's \texttt{produce} step runs each evolved generator for a certain amount of time, generating many seed inputs. As mentioned in Section~\ref{sec:gen-to-input}, a single generator can emit an unbounded number of inputs, so ELFuzz naturally will have a much larger raw corpus than \toolname{}, which has no \texttt{produce} step, and whose evolved pool \textit{is} the corpus. Also, the pool of inputs in \toolname{} will only grow if new inputs improve the coverage. This explains the massive disparity in corpus size before minimisation.

After minimisation with \texttt{afl-cmin}, seed corpora for all running configurations reduced in size drastically. \toolname{}'s jsoncpp corpus shrinks from 116 to just 3 inputs (2.6\%), showing that the directly generated inputs are highly redundant for coverage. Upon further inspection of the inputs, \textit{none} of them were valid JSON, meaning none of the inputs were able to bypass the parser of the jsoncpp library. This indicates that the model used is very poorly conditioned on producing valid JSON, and does not ``understand'' the syntax. \toolname{}'s cpython3 corpus shrinks from 1,264 to 118 inputs (9.3\%), better than jsoncpp, but still indicates a high level of redundancy. Only about a third of the inputs in the cpython3 corpus were valid Python, so the model is better at producing Python. See Section~\ref{sec:eval-validity} for further discussion of input validity.

Comparing ELFuzz to \toolname{}, the former produces a much larger corpus, even after minimisation. This would indicate that the model is much better at producing generators that can emit valid inputs, than it is at producing the inputs directly, which challenges our hypothesis. The generator abstraction appears to shield ELFuzz from the model's weakness in directly generating valid inputs, since a generator can encode the structure of the input format once and emit many valid instances, whereas direct generation must get each input right on its own. This is a significant factor in the seed-coverage gap observed in RQ1.

\begin{table}[h]
    \centering
    \begin{tabular}{lrrrr}
        \toprule
        & \multicolumn{2}{c}{\textbf{ELFuzz}} & \multicolumn{2}{c}{\textbf{\toolname{}}} \\
        \cmidrule(lr){2-3}\cmidrule(lr){4-5}
        \textbf{SUT} & \textbf{Before} & \textbf{After} & \textbf{Before} & \textbf{After} \\
        \midrule
        jsoncpp  & 2{,}800{,}000 & 205 & 116     & 3   \\
        cpython3 & 880{,}000     & 822 & 1{,}264 & 118 \\
        \bottomrule
    \end{tabular}
    \caption{Seed-corpus size before and after minimisation with \texttt{afl-cmin}. ``Before'' is the raw corpus handed to \texttt{afl-cmin}: for ELFuzz the output of its 200-second \texttt{produce} step, for \toolname{} the evolved pool that doubles as its corpus. ``After'' is the minimised, coverage-preserving corpus that seeds the AFL++ campaign.}
    \label{tab:seedcount}
\end{table}

\subsection{Coverage during pool evolution (\toolname{})}
\label{sec:eval-evolution}

Figure~\ref{fig:direct-evolution} traces the coverage held by \toolname{}'s pool over the course of evolution, for both SUTs. We show \toolname{} alone here, as its evolved pool is exactly the corpus it hands to AFL++, so the curve is a faithful picture of what the search produces. The equivalent curve for ELFuzz would not be comparable, because ELFuzz's per-generation coverage reflects the coverage achieved by a small number of generated inputs during evolution. This is separate from the full dedicated production step after synthesis, which produces many more inputs per generator, thereby resulting in much different coverage.

These use the synthesis-phase (\texttt{getcov}) instrumentation, so they are not on the same scale as the \texttt{afl-showmap} seed-coverage figures reported in Section~\ref{sec:eval-rq1}.

\begin{figure}[H]
    \centering
    \begin{minipage}[t]{0.49\linewidth}\centering
        \directcovtimeseries{jsoncpp}{evolution}{gen_index}{union_edges}{Generation}{Pool coverage (edges)}\\[2pt]
        {\small (a) jsoncpp}
    \end{minipage}\hfill
    \begin{minipage}[t]{0.49\linewidth}\centering
        \directcovtimeseries{cpython3}{evolution}{gen_index}{union_edges}{Generation}{Pool coverage (edges)}\\[2pt]
        {\small (b) cpython3}
    \end{minipage}
    \caption{Coverage held by \toolname{}'s pool across generations (the running union of edges over all retained inputs, on the \texttt{getcov} instrumentation). Generation $-1$ is the pre-generated initial seed set. Each later generation evaluates a batch of $1{,}000$ candidates.}
    \label{fig:direct-evolution}
\end{figure}

Both curves share the same overall shape: a steep initial climb followed by a long, decelerating tail. The first mutated generation alone accounts for most of the gain - on jsoncpp the pool jumps from $330$ edges in the pre-generated seed set to $545$ at generation~$0$, and on cpython3 from $16{,}427$ to $21{,}889$ - after which each generation contributes progressively less. This is the expected signature of coverage-guided search: early mutations cheaply reach the easily-accessible parts of the SUT, while later ones must work harder for what remains.

jsoncpp plateaus much sooner and more completely than cpython3. On jsoncpp the curve flattens almost completely after generation~$7$ at $592$ edges. It then rises to only $625$ by generation~$69$, a gain of about $30$ edges spread across the final sixty generations. However on cpython3, it continues to rise, but at a decelerating rate, until the end of the budget.
At generation~$56$ it has reached $25{,}225$ edges, and it is still rising at generation~$60$ at $25{,}261$ edges. At roughly $25$k edges it has reached only a fraction of the interpreter's far larger reachable space, so a longer evolution budget would plausibly keep yielding new coverage, albeit at a diminishing rate. This is consistent with the difference in complexity between the two SUTs: jsoncpp is a much smaller simpler SUT as compared to cpython3, so its reachable frontier is smaller and easier to find, and the search saturates it sooner. A more complex SUT seems to be better suited to the evolutionary search of \toolname{}, as there is less chance of saturating the reachable frontier within a realistic budget, and more room for the search to find new coverage over time.

\subsection{Seed coverage (RQ1)}
\label{sec:eval-rq1}

Table~\ref{tab:seedcov} reports the coverage achieved by each tool's minimised seed corpus on its own, measured with \texttt{afl-showmap}. This is the most direct measurement of seed quality.

\begin{table}[h]
    \centering
    \begin{tabular}{lrr}
        \toprule
        \textbf{SUT} & \textbf{ELFuzz} & \textbf{\toolname{}} \\
        \midrule
        jsoncpp  & 568    & 163   \\
        cpython3 & 19{,}634 & 9{,}790 \\
        \bottomrule
    \end{tabular}
    \caption{Seed-corpus coverage (edges covered by the minimised corpus, before any mutation fuzzing) for each tool, measured with \texttt{afl-showmap}. Higher is better.}
    \label{tab:seedcov}
\end{table}

On both SUTs the generator-synthesising pipeline produces a higher-coverage seed corpus than direct input generation: ELFuzz leads by roughly $3.5\times$ on jsoncpp (568 vs 163 edges) and $2\times$ on cpython3 (19{,}634 vs 9{,}790). This difference is partly due to the way in which the two seed corpora are generated. ELFuzz's \texttt{produce} step generates a much larger raw corpus than \toolname{}, and \toolname{}'s seed corpus can only grow when new coverage is achieved. While improvements to coverage are part of the generator evolution process in ELFuzz, a single generator can achieve more coverage than was measured during evolution. With 880,000 / 2.8M inputs, ELFuzz's seed corpus is allowed to risk some redundancy in exchange for sheer quantity. Also, as mentioned in Section~\ref{sec:eval-seedcount}, the model used is much better at producing generators that can emit valid inputs, than it is at producing the inputs directly, allowing ELFuzz's seeds to reach past the initial parsing stage of the SUTs and cover more edges.

The difference in coverage is not just in quantity, as ELFuzz's seed corpora are almost entirely \textit{supersets} of our seed corpora. Table~\ref{tab:edgesets} shows the individual edges covered by each tool, and their relationship to each other.

\begin{table}[h]
    \centering
    \begin{tabular}{lrrrrr}
        \toprule
        \textbf{SUT} & \textbf{ELFuzz} & \textbf{\toolname{}} & \textbf{Shared} & \textbf{ELFuzz-only} & \textbf{\toolname{}-only} \\
        \midrule
        jsoncpp  & 568      & 163     & 163     & 405     & 0  \\
        cpython3 & 19{,}634 & 9{,}790 & 9{,}755 & 9{,}879 & 35 \\
        \bottomrule
    \end{tabular}
    \caption{Edges covered by each tool's minimised seed corpus and their overlap, measured with \texttt{afl-showmap} (so both tools sit in the same edge space). ``Shared'' is the intersection of the two edge sets; ``ELFuzz-only'' and ``\toolname{}-only'' are the edges unique to each tool. On both SUTs \toolname{}'s edges are almost entirely a subset of ELFuzz's.}
    \label{tab:edgesets}
\end{table}

On jsoncpp every one of \toolname{}'s 163 edges is also covered by ELFuzz, which reaches 405 more; \toolname{} contributes \textit{no} edge that ELFuzz misses. On cpython3, \toolname{} only achieves 35 unique edges not covered by ELFuzz, but ELFuzz covers 9{,}879 edges absent from \toolname{}'s corpus. Directly generated inputs are therefore exploring a subset of the generator-generated inputs, rather than a complementary part of the target. This is the strongest evidence against our hypothesis that direct generation would perform superior to indirect generation.

\subsection{Downstream AFL++ coverage (RQ2)}
\label{sec:eval-rq2}

\begin{figure}[H]
    \centering
    \begin{minipage}[t]{0.32\linewidth}\centering
        \covtimeseries[xmax=0.4,scale only axis,width=3cm,height=3cm]{jsoncpp}{afl}{time_h}{edges_found}{Time (hours)}{Edges found (AFL++)}\\[2pt]
        {\small (a.1) jsoncpp (first 0.4\,h)}
    \end{minipage}\hfill
    \begin{minipage}[t]{0.32\linewidth}\centering
        \covtimeseries[scale only axis,width=3cm,height=3cm]{jsoncpp}{afl}{time_h}{edges_found}{Time (hours)}{Edges found (AFL++)}\\[2pt]
        {\small (a.2) jsoncpp (full 6\,h)}
    \end{minipage}\hfill
    \begin{minipage}[t]{0.32\linewidth}\centering
        \covtimeseries[scale only axis,width=3cm,height=3cm]{cpython3}{afl}{time_h}{edges_found}{Time (hours)}{Edges found (AFL++)}\\[2pt]
        {\small (b) cpython3}
    \end{minipage}
    \caption{AFL++ coverage over time, seeded with each tool's minimised corpus; campaigns run for 6\,h. Panel (a.1) shows jsoncpp truncated to the first 0.4\,h, where both tools reach the 667-edge reachable frontier within the first $\sim$20 minutes and stay flat thereafter; (a.2) shows the same jsoncpp campaign over the full 6\,h; (b) shows cpython3 over the full 6\,h, where neither tool saturates.}
    \label{fig:afl}
\end{figure}

\begin{table}[h]
    \centering
    \begin{tabular}{llrrrrr}
        \toprule
        \textbf{SUT} & \textbf{Tool} & \textbf{Edges} & \textbf{Bitmap} & \textbf{Corpus} & \textbf{Execs} & \textbf{Crashes} \\
        \midrule
        \multirow{2}{*}{jsoncpp}
          & ELFuzz     & 667 & 12.74\% & 1311 & 166.9\,M & 0 \\
          & \toolname{} & 667 & 12.74\% & 1468 & 141.7\,M & 0 \\
        \midrule
        \multirow{2}{*}{cpython3}
          & ELFuzz     & 24{,}277 & 18.20\% & 6043 & 3.31\,M & 0 \\
          & \toolname{} & 22{,}523 & 16.88\% & 6636 & 3.33\,M & 0 \\
        \bottomrule
    \end{tabular}
    \caption{End-of-campaign AFL++ statistics (one repetition): edges found, bitmap coverage, final corpus size, total executions, and saved crashes. cpython3's far lower execution rate (millions of execs, vs hundreds of millions for jsoncpp) reflects how much heavier each interpreter run is.}
    \label{tab:aflstats}
\end{table}

The two SUTs behave very differently. On jsoncpp the campaigns converge to the \textit{identical} final coverage - 667 edges, 12.74\% of the map - even though the ELFuzz seed corpus started out covering $3.5\times$ as many edges as the \toolname{} one. AFL++ saturates jsoncpp's reachable frontier within the first few minutes (both curves are flat for almost the entire campaign), so the large seed-quality difference is erased entirely; it affects only how quickly the plateau is reached, not where it lands. On cpython3, by contrast, neither campaign comes close to saturation. They finish at 18.20\% and 16.88\% of a much larger map, and are both still climbing. Here the seed advantage persists: ELFuzz stays ahead throughout and ends with 24{,}277 edges compared to 22,523 edges for \toolname{}. The campaign narrows the relative gap (from $2\times$ at the seed stage to about $8\%$ at the end) but does not close it.

\subsection{Validity of the generated inputs}
\label{sec:eval-validity}

Of the 116 inputs in \toolname{}'s jsoncpp corpus, 0 were valid JSON. Upon inspection, apart from simple syntax errors, many of the inputs contained HTML tags, or JavaScript code, either directly in the input, in string literals, or after a \texttt{<file\_sep>} token - this token is used by the model to indicate a separation between files \cite{hui_2024_qwen25coder}. Clearly the model was not trained on raw JSON files, as they were often interpreted as part of a HTML or JavaScript context. This is a significant weakness for direct generation, as it means the model is not able to leverage its knowledge of the input format to produce valid inputs that can reach deeper into the SUT.

In contrast, 422 out of the 1,264 inputs in \toolname{}'s cpython3 corpus were valid Python. The model is much better at producing valid Python than valid JSON, which also can explain the good performance of ELFuzz, evolving python programs too. This allowed \toolname{} to produce a significant number of valid inputs that could reach deeper into the SUT, which is reflected in the higher seed coverage compared to jsoncpp. Still, only 33.4\% of the seed corpus was valid Python, again severely limiting where in the SUT the inputs could reach, and therefore the coverage they could achieve. This is a major factor in the seed-coverage gap between ELFuzz and \toolname{} on cpython3, as ELFuzz's generator abstraction allows it to produce many valid inputs even if the model is not good at directly generating them.

\section{Discussion}
\label{sec:eval-discussion}

Here are some key takeaways from the results.

Using LLMs to generate seed inputs for fuzzing is less valuable for simpler SUTs, as demonstrated by jsoncpp. Even from 3 inputs after minimisation, the SUT was saturated within $\sim$20 minutes. For cpython3, a much more complex SUT, neither ELFuzz nor \toolname{} saturated within the 6-hour budget. It is unlikely that cpython3 would saturate with an extended time budget (e.g. 24 or 48 hours), so there is more benefit to be had in using LLMs for fuzzing richer SUTs, such as other programming languages.

For cpython3, the quality of the seed corpus matters for the final coverage, as demonstrated by the persistent gap between ELFuzz and \toolname{}. AFL++ cannot exhaust the target within the budget, so it amplifies the seed corpus and preserves its quality difference throughout. This suggests that for realistic, large targets where saturation is unlikely, investing in a high-quality seed corpus is important for achieving better coverage.

The model's ability to produce valid inputs is a critical factor for direct generation. \toolname{}'s jsoncpp corpus contained no valid JSON, which severely limited its coverage and AFL++'s ability to explore the SUT. In contrast, \toolname{}'s cpython3 corpus had a higher number of valid Python inputs, but still a low percentage overall. The generator abstraction of ELFuzz allows it to bypass the model's weakness in generating many valid python programs, as a single generator can encode the structure of the input format and emit many valid instances.

With the evidence of these two SUTs, the directness hypothesis is not yet supported.

\section{Threats to Validity}
\label{sec:eval-threats}

These are early results - a single repetition, a small local model, and only two SUTs - and a fuller comparison across more SUTs and repetitions is needed before drawing firm conclusions.

\textbf{Internal validity.} All results in this chapter come from a single repetition per configuration, so run-to-run variance is not measured. In addition, while both tools use the same model for their evolution loop, \toolname{}'s initial pool is pre-generated by a different, more capable model, Claude Opus 4.8 (Section~\ref{sec:init}), so some of the coverage held by its pool is effectively free; ELFuzz begins its own evolution from a naive seed fuzzer that emits purely random bytes, and does not have this initial coverage advantage.

\textbf{Construct validity.} The pool-evolution curves of Section~\ref{sec:eval-evolution} use the synthesis-phase (\texttt{getcov}) instrumentation, so they are not on the same scale as the \texttt{afl-showmap} figures reported in Sections~\ref{sec:eval-rq1} and~\ref{sec:eval-rq2}.

\textbf{External validity.} Hardware constraints limited us to a small local model, Qwen2.5-Coder-1.5B, where the original ELFuzz experiments used the far larger CodeLlama-13B \cite{chen_2025_elfuzz}, and our campaigns ran for 6 hours rather than the 24 used in the original evaluation. With only two SUTs, jsoncpp and cpython3, the results may not generalise to other input formats and richer targets.

